% Sewall Wright, "Statistical genetics and evolution", Bull. Amer. Math. Soc. 48 (1942), first page, reset in Mills 8A.
\documentclass[11pt]{article}
\usepackage{bulletin1947}
\begin{document}
\setcounter{page}{223}
\received{The sixteenth Josiah Willard Gibbs Lecture, delivered at Chicago, Illinois, September 3, 1941, under the auspices of the American Mathematical Society; received by the editors October 20, 1941.}
\articletitle{STATISTICAL GENETICS AND EVOLUTION}{sewall wright}
\head{Introduction.} When Darwin developed the theory of evolution by natural selection, practically nothing was known of hereditary differences beyond their existence. Since 1900, a body of knowledge on the mechanism of heredity and on mutation has been built up by experiment that challenges any field in the biological sciences in the extent and precision of its results. The implications for evolution are not, however, immediately obvious. It is necessary to work out the statistical consequences.

Studies in the field of statistical genetics began shortly after the rediscovery of Mendelian heredity in 1900. Those of J. B. S. Haldane [\textbf{7}] and R. A. Fisher [\textbf{4}] have been especially important with respect to the application to evolution. My own approach to the subject came through experimental studies conducted in the U. S. Bureau of Animal Industry on the effects of inbreeding, crossbreeding and selection on populations of guinea pigs [\textbf{21, 22, 23, 37}] and through the attempt to formulate principles applicable to livestock breeding [\textbf{19, 20, 24, 25, 13, 34}]. On moving into the more academic atmosphere of the University of Chicago, I have become more directly concerned with the problem of evolution.

I should note that the deductive approach, to which I shall confine myself here, involves many questions that can only be settled by observation and experimental work on natural populations and that a remarkable resurgence of interest in such work is in progress [\textbf{2, 9}].

\medskip
\head{Postulates.} It will be desirable to begin with a brief review of the more important factors of which account must be taken.

The basic fact of modern genetics is that heredity can be analyzed into separable units, ``genes,'' whose most essential property is that of duplicating themselves with extraordinary precision, irrespective of the characteristics of the organism in whose cells they are carried. We shall restrict consideration to changes in the system of genes and aggregates of genes (chromosomes). There are relatively rare and obscure hereditary changes which must be attributed to other cell components but our knowledge of these does not warrant the elaboration of a statistical theory.

Fortunately the same theory applies to a large extent to gene muta
\end{document}
